% !TeX root = ../../../../main.tex
\subsubsection{Clinician Annotation Quality Review}
\label{ch06:apd:clinician_quality_review}

\ChSixHeadingOnly
\paragraph{Review protocol}

In addition to the quantitative metrics used in the main evaluation, a physician reviewed extracted timelines for the \texttt{sepsis-10} subset.
For each case report, the clinician assigned: (i) a rank ordering (1 = best, 7 = worst) of the annotations, consisting of the manual annotation and six LLM-generated annotations, and (ii) a categorical assessment (Excellent, Good, Acceptable, or Poor) based on extraction completeness, temporal ordering accuracy, and timing plausibility.
This review provides a qualitative assessment of annotation style and clinical usability, complementing the quantitative event-matching and temporal-alignment metrics reported in the main text.
The reviewed annotation bundle included O1-preview as the OpenAI annotation; replacing GPT-5 which was evaluated separately in the main quantitative results in Table~\ref{ch06:tab:results}.

The clinician review results are reported in Table~\ref{ch06:tab:sep10review}.

\ChSixParagraph{Events requiring temporal reasoning or composition.}
Example: ``passed away'' from ``he was re-admitted to the medical ICU for severe sepsis and multiorgan failure and passed away around 6 months after his initial diagnosis with NHL, despite maintaining a remission status.''
This requires a two-step inference: (1) timing of NHL diagnosis, and (2) adding ``around 6 months'' to that time.
While this is a simple example, more complicated chains of relative timing were noted, sometimes involving other events whose timings are themselves uncertain.
As black-box methods, there is no explicit probabilistic specification for these chains, and empirically we see variable responses.
For this case, the manual annotation is 4383 hours (about 6 months), O1-preview gives 4368 hours (about 6 months), DeepSeek-R1 gives 4032 hours (about 6 months), and Llama 3.3 70B gives 1440 hours (about 2 months); compared to other events the temporal ordering is correct, but the relative timing is off by months.
